\section{Related Work}
\label{sec:related}

\paragraph{Image cropping.}
Automatic image cropping has evolved from rule-based composition
heuristics~\cite{cgs2017} through reinforcement
learning~\cite{a2rl2018} to fully supervised approaches trained on
large-scale human annotations~\cite{gaic2019, gaicd2020}.  While
effective, these methods require task-specific training and do not
generalize easily across datasets or cropping objectives.
\baseline~\cite{cropper2025} circumvents this limitation entirely by
using VLM in-context learning, achieving competitive results without any
gradient updates.

\paragraph{Vision-Language Models and ICL.}
Large VLMs such as Gemini~\cite{gemini2024},
Mantis~\cite{mantis2024}, and Idefics2~\cite{idefics2} have
demonstrated strong multi-image understanding through interleaved
image--text prompting.  In-context learning---providing task
demonstrations at inference time---has become a standard paradigm for
adapting these models to new tasks without fine-tuning.  \baseline
exploits this by retrieving CLIP-similar training images with annotated
crops and presenting them as ICL demonstrations.

\paragraph{Aesthetic scoring.}
VILA~\cite{vila2024} learns image aesthetics from user comments via
vision-language pretraining and provides a scalar aesthetic quality
score.  \baseline combines VILA with a heuristic Area scorer to evaluate
crop candidates.  The LAION aesthetic predictor~\cite{laion_aes2022}
offers an alternative lightweight aesthetic score.  In this work, we
examine the interaction between scorer design and iterative refinement,
showing that the Area component introduces a systematic bias that
undermines the refinement loop.
